%% toc: P3 26 26.5,26.6,26.7
\section{Proofs: Vaught's Section 6 routes and the Vaught ordinal}
\label{sec:annex-p03-w03-p03-016-vaught}

Throughout, languages are countable, and the theories \(T_{r+2}\) and
\(T'_{r+2}\) are those of \cref{def:p03-vaught-colored-orders}. The
compactness theorem is \cref{thm:fol-compactness}, the downward
L\"owenheim--Skolem theorem is \cref{thm:downward-ls}, and the union of an
elementary chain is handled by \cref{thm:elementary-chain}. A theory is
\emph{complete} when it is consistent and, for every sentence \(\varphi\)
of its language, proves \(\varphi\) or \(\neg\varphi\).

\begin{annexitem}{annex:p03-w03-p03-016-source-routes}{thm:p03-w03-p03-016-source-routes}{Completeness by Vaught's two routes}
\begin{proof}
\emph{A completeness criterion.} Let \(Q\) be a consistent theory in a
countable language all of whose models are infinite, and suppose any two
denumerable models of \(Q\) are elementarily equivalent. Then \(Q\) is
complete. Indeed, if \(Q\) proved neither \(\varphi\) nor \(\neg\varphi\),
the completeness theorem would give models \(M\entails Q\cup\{\neg\varphi\}\)
and \(N\entails Q\cup\{\varphi\}\). Downward L\"owenheim--Skolem gives
countable elementary substructures \(M'\preccurlyeq M\) and
\(N'\preccurlyeq N\). They are models of \(Q\), hence infinite, hence
denumerable, and they disagree on \(\varphi\). Both theories below are
consistent, because the annex proof of
\cref{thm:p03-vaught-finite-counts} constructs denumerable models in every
case.

\emph{Clause 1.} Let \(M\) and \(N\) be denumerable models of
\(T_{r+2}\), and fix \(k\). In \(M\) the points
\(c_0^M<\cdots<c_k^M\) cut the domain into the open pieces
\[
 I_{-1}=\{x:x<c_0\},\quad I_j=\{x:c_j<x<c_{j+1}\}\ (j<k),\quad
 I_k=\{x:c_k<x\},
\]
and likewise in \(N\). Each piece is nonempty, because the order is dense
and has no endpoints. With the induced order it is dense without endpoints:
a point of the piece has neighbours on both sides inside it, by density
between it and the cutting points and by the absence of endpoints. Each
color \(U_t\) is dense in the piece, because it is dense in the whole
order. Each piece is countably infinite. By Skolem's fact as Vaught uses it,
proved by the colored back-and-forth argument at the start of the annex
proof of \cref{thm:p03-vaught-finite-counts}, there is a color-preserving
order isomorphism \(f_j:I_j^M\to I_j^N\) for each \(j\) from \(-1\) to
\(k\). Let \(f\) be the union of these maps together with
\(c_i^M\mapsto c_i^N\) for \(i\le k\). The pieces and the cutting points
partition both domains, so \(f\) is a bijection. It preserves the order
inside each piece by construction and between pieces and cutting points
because both sides list them in the same order. It preserves colors on the
pieces by construction and at the cutting points because every \(c_i\) lies
in \(U_0\). Hence \(f\) is an isomorphism of the reducts to
\(<,U_0,\ldots,U_{r-1},c_0,\ldots,c_k\). A sentence of the full language
mentions finitely many constants, so it lies in such a reduct language for
some \(k\), and \(M\) and \(N\) agree on it. Every model of \(T_{r+2}\) is
infinite, being a dense order without endpoints. The completeness criterion
applies.

\emph{Notation for clause 2.} In a model of \(T'_{r+2}\) write \([x]\)
for the \(R\)-class of \(x\). Since the order is constant on classes and
the quotient is a dense linear order, \(x<y\) holds exactly when
\([x]\) precedes \([y]\) in the quotient. In particular \(<\) is
transitive, and \(x<y\) implies \([x]\ne[y]\). For a positive integer
\(i\) let \(\mathrm{Size}_i(y)\) be the formula
\[
 \exists y_1\cdots\exists y_i\Bigl(\bigwedge_{a<b\le i}y_a\ne y_b
 \wedge\bigwedge_{a\le i}R(y,y_a)
 \wedge\forall z\bigl(R(y,z)\rightarrow\bigvee_{a\le i}z=y_a\bigr)\Bigr),
\]
and let \(\beta_i(x)\) be \(\forall y(\mathrm{Size}_i(y)\rightarrow y<x)\).
Put \(\Gamma(x)=\{\beta_i(x):i\ge1\}\). In a model, \(C_i\) is the unique
class of size \(i\), so an element \(e\) satisfies every member of
\(\Gamma\) exactly when \([e]\) is an upper bound of all the finite classes.
The axiom relating the classes of sizes \(i\) and \(i+1\), with
transitivity, gives \(C_i<C_{i+1}<\cdots\); consequently every element of
\(C_{N+1}\) satisfies \(\beta_1,\ldots,\beta_N\). By the classification in
the annex proof of \cref{thm:p03-vaught-finite-counts}, which uses only the
axioms, a denumerable model falls into exactly one of the following cases,
and two denumerable models in the same case are isomorphic: \(\Gamma\) is
not realized (the finite classes are cofinal); \(\Gamma\) is realized but
the classes realizing it have no least member; or there is a least class
realizing \(\Gamma\), of color \(U_t\) for one of the \(r\) values of
\(t\). Every case occurs.

\emph{Step 1: the cofinal case.} Let \(A\) be a denumerable model in which
the finite classes are cofinal. Add a constant \(c\) and consider
\(\operatorname{Th}(A)\cup\{\beta_i(c):i\ge1\}\). A finite subset mentions
\(\beta_i(c)\) only for \(i\le N\), and it holds in \(A\) when \(c\) is
interpreted as any element of \(C_{N+1}^A\). By compactness the whole set
has a model, and downward L\"owenheim--Skolem gives a countable elementary
substructure \((A'',c'')\) of it. It is a model of
\(\operatorname{Th}(A)\), so \(A''\equiv A\) and \(A''\) is a model of
\(T'_{r+2}\). It contains infinitely many finite classes, so it is
denumerable. The element \(c''\) realizes \(\Gamma\), so the finite classes
of \(A''\) are bounded: \(A''\) is in the second or the third case.

\emph{Step 2: the least-upper-bound case.} Let \(B_0\) be a denumerable
model with a least class realizing \(\Gamma\), and let \(\ell\) be an
element of it. We build denumerable models
\(B_0\preccurlyeq B_1\preccurlyeq\cdots\). Given \(B_k\), let \(E_k\) be
the set of elements of \(B_k\) satisfying every \(\beta_i\) in \(B_k\).
Expand the language by a constant \(\underline b\) for each \(b\in B_k\)
and one further constant \(d\), and let \(\Sigma_k\) consist of the
elementary diagram of \(B_k\), all \(\beta_i(d)\), and all
\(d<\underline e\) for \(e\in E_k\). A finite subset of \(\Sigma_k\)
mentions \(\beta_i(d)\) only for \(i\le N\) and finitely many
\(e_1,\ldots,e_m\in E_k\). Interpret every \(\underline b\) as \(b\) and
\(d\) as an element of \(C_{N+1}^{B_k}\). The diagram sentences hold, the
\(\beta_i(d)\) with \(i\le N\) hold by the preceding paragraph, and
\(d<e_j\) holds because \(e_j\) satisfies \(\beta_{N+1}\). By compactness
\(\Sigma_k\) has a model \(D_k\). The map sending \(b\) to the value of
\(\underline b\) is an elementary embedding of \(B_k\), so after renaming
elements we may take \(B_k\preccurlyeq D_k\). Downward
L\"owenheim--Skolem gives a countable \(B_{k+1}\preccurlyeq D_k\)
containing \(B_k\) and the value \(d_k\) of \(d\). For a formula with
parameters from \(B_k\), truth in \(B_k\), in \(D_k\) and in \(B_{k+1}\)
coincide, so \(B_k\preccurlyeq B_{k+1}\). The first-order facts
\(\beta_i(d_k)\) and \(d_k<e\) for \(e\in E_k\) pass from \(D_k\) to
\(B_{k+1}\) by elementarity. Since \(B_{k+1}\) contains \(B_k\), it is
denumerable.

Let \(B_\omega\) be the union of the chain. It is countable, and
\(B_k\preccurlyeq B_\omega\) for every \(k\) by the elementary chain
theorem. In particular \(B_\omega\equiv B_0\), so \(B_\omega\) is a
denumerable model of \(T'_{r+2}\). The element \(\ell\) satisfies every
\(\beta_i\) in \(B_0\), hence in \(B_\omega\), so the finite classes of
\(B_\omega\) are bounded. Suppose some class of \(B_\omega\) were least
among the classes realizing \(\Gamma\), and let \(e\) be a member. Then
\(e\in B_k\) for some \(k\), and by elementarity \(e\) satisfies every
\(\beta_i\) in \(B_k\), so \(e\in E_k\). The element \(d_k\in B_{k+1}\)
satisfies every \(\beta_i\) and \(d_k<e\) in \(B_{k+1}\), hence in
\(B_\omega\). Thus \([d_k]\) realizes \(\Gamma\) and precedes \([e]\), a
contradiction. Hence \(B_\omega\) is in the second case and is elementarily
equivalent to \(B_0\).

\emph{Conclusion of clause 2.} Fix a denumerable model \(D\) in the second
case. Every model in the third case is elementarily equivalent, by Step~2,
to a model in the second case, which is isomorphic to \(D\). The cofinal
model \(A\) is elementarily equivalent, by Step~1, to a model in the second
or third case, hence to \(D\). Thus all denumerable models of
\(T'_{r+2}\) are elementarily equivalent. Every model of \(T'_{r+2}\) is
infinite, since its classes \(C_1,C_2,\ldots\) are pairwise disjoint and
nonempty. The completeness criterion applies.
\end{proof}
\end{annexitem}

\begin{annexitem}{annex:p03-w03-p03-016-theorem-6-1}{thm:p03-w03-p03-016-theorem-6-1}{Theorem 6.1 over all natural numbers}
\begin{proof}
\emph{A finite model fixes the size.} If a complete theory \(T\) has a
model of finite size \(d\), let \(\varepsilon_d\) be the first-order sentence
saying that there are exactly \(d\) elements. It holds in that model, so
\(T\) does not prove \(\neg\varepsilon_d\); by completeness \(T\) proves
\(\varepsilon_d\), and every model of \(T\) has exactly \(d\) elements.
In particular \(T\) has no denumerable model.

\emph{The case \(n=0\).} Let the language have one unary relation symbol
\(P\), and let \(A\) have a single element, with \(P^A\) the whole domain.
The set \(\operatorname{Th}(A)\) of sentences true in \(A\) is consistent,
having \(A\) as a model, and contains \(\varphi\) or \(\neg\varphi\) for
every sentence \(\varphi\); so it is complete. By the preceding paragraph
it has no denumerable model, so it has exactly \(0\) of them.

\emph{The cases \(n=1\) and \(n\ge3\).} By
\cref{thm:p03-vaught-finite-counts}, dense linear order without endpoints
has exactly one denumerable model, and for \(n=r+2\) with \(r=n-2\ge1\)
the complete theory \(T_{r+2}\) has exactly \(n\). Their languages consist
of relation symbols and individual constants.

\emph{The case \(n=2\) is impossible.} Let \(T\) be complete in a countable
language. If \(T\) has a finite model, it has no denumerable model, so its
count is \(0\). Otherwise all its models are infinite, so its countable
models are exactly its denumerable models, and their number
\(I(T,\aleph_0)\) is not \(2\) by
\cref{thm:p03-vaught-never-two-original}.

\emph{An infinite model gives a denumerable one.} Let \(T\) be complete in
a countable language with an infinite model \(M\). Downward
L\"owenheim--Skolem, applied with the empty set of required elements, gives
\(N\preccurlyeq M\) with at most \(\aleph_0\) elements. For each \(k\) the
sentence saying that there are at least \(k\) elements holds in \(M\), hence
in \(N\). So \(N\) is infinite, hence denumerable, and it is a model of
\(T\).
\end{proof}
\end{annexitem}

\begin{annexitem}{annex:p03-w03-p03-016-vo-cardinal}{thm:p03-w03-p03-016-vo-cardinal}{A countable Vaught ordinal gives the cardinal dichotomy}
\begin{proof}
\emph{Isomorphism invariance.} If \(h:M\to N\) is an isomorphism, then for
every \(L_{\omega_1\omega}\)-formula \(\psi(\bar x)\) and tuple \(\bar a\)
from \(M\), \(M\entails\psi(\bar a)\) exactly when
\(N\entails\psi(h\bar a)\). This is proved by induction on \(\psi\):
atomic formulas are preserved because \(h\) is an isomorphism; negation and
countable conjunctions and disjunctions are handled pointwise; and for an
existential quantifier a witness \(b\) in \(M\) gives the witness \(hb\) in
\(N\), while a witness \(b'\) in \(N\) gives \(h^{-1}b'\) in \(M\).
Consequently isomorphic models realize the same
\(\Pi^{\mathrm{in}}_\beta\)-types.

\emph{Clause 1.} Let \(\beta=\operatorname{vo}(\sigma)<\omega_1\). By
definition one of the two alternatives holds at \(\beta\).

Suppose that there are only countably many models of \(\sigma\), with Scott
ranks below \(\beta\). Then \(I(\sigma,\aleph_0)\le\aleph_0\).

Suppose instead that continuum many \(\Pi^{\mathrm{in}}_\beta\)-types are
realized among the countable models of \(\sigma\). Let
\(\kappa=I(\sigma,\aleph_0)\) and choose representatives \(M_\xi\)
(\(\xi<\kappa\)) of the isomorphism types. By isomorphism invariance every
realized type is realized by a tuple of some \(M_\xi\). Each \(M_\xi\) is
countable and so has countably many finite tuples. Hence the set of
realized types has cardinality at most \(\kappa\cdot\aleph_0\). If
\(\kappa\le\aleph_0\), this is at most \(\aleph_0<2^{\aleph_0}\),
contradicting the hypothesis. So \(\kappa\) is infinite,
\(\kappa\cdot\aleph_0=\kappa\), and \(\kappa\ge2^{\aleph_0}\). On the other
hand, a countable language has at most \(2^{\aleph_0}\) countable
isomorphism types (\cref{thm:p03-vaught-cardinal-boundaries}). Therefore
\(I(\sigma,\aleph_0)=2^{\aleph_0}\). In either case the dichotomy holds.

\emph{Clause 2.} If \(\alpha<\omega_1\), then \(\alpha+\omega\) is the
union of \(\alpha\) and the countable set
\(\{\alpha+m:m<\omega\}\), so it is a countable ordinal and
\(\operatorname{vo}(\sigma)\le\alpha+\omega<\omega_1\). The
\(\omega\)-Vaught conjecture, as formulated in
\cref{def:p03-vaught-scott-comparison}, includes the assertion that every
sentence has a countable Vaught ordinal, so by clause~1 it implies the
infinitary cardinal dichotomy for every sentence.

\emph{Clause 3.} The countable models of \(\sigma\) form a class of
countable structures in a countable language. Under the continuum
hypothesis \cref{thm:p03-vaught-cardinal-boundaries} gives the bare
cardinal dichotomy for every such class, so every \(\sigma\) satisfies it.
The statement that every \(\sigma\) satisfying the dichotomy has a countable
Vaught ordinal therefore has a hypothesis true of every \(\sigma\), and it
is equivalent to the statement that every \(\sigma\) has a countable Vaught
ordinal.

\emph{Verification of \cref{def:p03-w03-p03-016-independent-predicates}.}
In the structure \(M_0\) of finite subsets of \(\omega\), the sentence
\(\theta_{A,B}\) is witnessed by \(x=A\): every \(i\in A\) belongs to
\(A\), and no \(j\in B\) does, since \(A\) and \(B\) are disjoint. So
\(M_0\entails\sigma_P\), and \(M_0\) is countable. Fix
\(S\subseteq\omega\) and a new constant \(c\). Consider the first-order set
consisting of all \(\theta_{A,B}\), all \(P_i(c)\) for \(i\in S\), and all
\(\neg P_j(c)\) for \(j\notin S\). A finite subset mentions only indices
below some \(N\), and it holds in \(M_0\) with \(c\) interpreted as the
finite set \(S\cap N\). By compactness the set has a model, and downward
L\"owenheim--Skolem gives a countable elementary substructure containing
the value of \(c\). That structure satisfies every \(\theta_{A,B}\), hence
the countable conjunction \(\sigma_P\), and the value of \(c\) lies in
\(P_i\) exactly for \(i\in S\). If \(S\ne S'\) and \(i\) lies in exactly
one of them, the corresponding elements disagree on the quantifier-free
formula \(P_i(x)\), so they realize different
\(\Pi^{\mathrm{in}}_0\)-types. This gives at least \(2^{\aleph_0}\)
realized \(\Pi^{\mathrm{in}}_0\)-types. There are only countably many
finitary quantifier-free formulas, so each type is a subset of a countable
set, and there are at most \(2^{\aleph_0}\) of them. Hence exactly
continuum many \(\Pi^{\mathrm{in}}_0\)-types are realized, the first
alternative holds at \(0\), and \(\operatorname{vo}(\sigma_P)=0\). Clause~1
gives \(I(\sigma_P,\aleph_0)=2^{\aleph_0}\).
\end{proof}
\end{annexitem}
